% !TEX root = ../DoAn.tex
\documentclass[../DoAn.tex]{subfiles}
\begin{document}

Chương~\ref{chapter:Experiment} đã trình bày thiết kế, xây dựng và đánh giá tổng thể hệ thống. Chương này tập trung vào những giải pháp mang tính đóng góp riêng của đồ án --- các phần do tác giả tự thiết kế thay vì chỉ sử dụng thư viện có sẵn. Mỗi đóng góp được trình bày theo ba phần: bài toán cần giải quyết, giải pháp đã triển khai và kết quả đạt được. Sáu đóng góp được trình bày gồm: chuỗi phân giải vị trí vệ tinh GNSS theo độ ưu tiên (Mục~\ref{section:5.1}); cơ chế chỉ đường trực tiếp bằng đo chuyển động thị giác khi GNSS suy giảm (Mục~\ref{section:5.live-routing-fusion}); cơ chế so sánh hiệu năng hai thuật toán Optical Flow (Mục~\ref{section:5.2}); quy trình tải lên và xử lý video RAFT bất đồng bộ trên máy chủ (Mục~\ref{section:5.3}); cơ chế bám đối tượng trong luồng Optical Flow camera trực tiếp (Mục~\ref{section:5.4}); và bộ dựng giao diện kính lỏng thời gian thực bằng AGSL kèm cơ chế suy biến mượt (Mục~\ref{section:5.liquidglass}).

\section{Chuỗi phân giải vị trí vệ tinh GNSS theo độ ưu tiên}
\label{section:5.1}
\subsection{Bài toán}
Trên Android, trạng thái vệ tinh hầu như luôn cung cấp azimuth và elevation, nhưng tọa độ vệ tinh trong không gian thì không phải lúc nào cũng có. Một số thiết bị cho phép trích xuất PVT chi tiết từ dữ liệu đo GNSS, song khả năng này không ổn định giữa các phiên bản Android và chipset. Nếu chỉ dựa vào azimuth/elevation, mô hình 3D vẫn hiển thị được vệ tinh ở vị trí xấp xỉ nhưng bỏ phí những nguồn dữ liệu chính xác hơn khi chúng sẵn có.

Bài toán đặt ra là xây dựng một cơ chế phân giải vị trí vệ tinh có thứ tự ưu tiên rõ ràng: luôn dùng nguồn tốt nhất hiện có, không để giao diện trống khi thiếu dữ liệu, và ghi lại nguồn của mỗi vị trí để cùng một điểm hiển thị vẫn phản ánh đúng mức độ tin cậy của nó.

\subsection{Giải pháp}
Đồ án triển khai một lớp theo dõi vệ tinh đóng vai trò điểm tập trung của toàn bộ dữ liệu vệ tinh. Lớp này duy trì ba vùng đệm: PVT mới nhất theo từng vệ tinh, broadcast ephemeris từ IGS và bản ghi quỹ đạo từ CelesTrak. Với mỗi vệ tinh quan sát được, hệ thống lần lượt thử các nguồn theo thứ tự ưu tiên giảm dần và dừng ngay khi có kết quả hợp lệ, nhờ vậy không buộc mọi thiết bị phải có cùng khả năng GNSS. Mỗi nguồn đi kèm điều kiện hợp lệ riêng: PVT chỉ dùng khi chưa quá hạn 10 giây; broadcast ephemeris chỉ dùng khi thời điểm bản tin gần thời điểm quan sát trong ngưỡng 12 giờ; CelesTrak dùng cho các chòm được hỗ trợ; cuối cùng là vị trí xấp xỉ tính từ azimuth, elevation và vị trí người quan sát. Bảng~\ref{table:gnss-source-priority} tóm tắt bốn mức ưu tiên này.

Điểm cốt lõi của giải pháp là mỗi kết quả phân giải không chỉ mang tọa độ mà còn mang nhãn nguồn vị trí và nguồn lịch quỹ đạo. Thông tin nguồn được chuyển tiếp tới bộ dựng hình và hộp thoại chi tiết, nhờ đó giao diện có thể giải thích dữ liệu đang hiển thị, còn nhật ký phát triển có thể phát hiện thời điểm hệ thống chuyển giữa PVT và các nguồn dự phòng.

\subsection{Kết quả đạt được}
Giải pháp giúp module GNSS hoạt động ổn định ở nhiều mức khả dụng dữ liệu khác nhau. Khi thiết bị hỗ trợ PVT, hệ thống ưu tiên tọa độ ECEF thật và tính được vận tốc quán tính; khi không có PVT, ứng dụng vẫn hiển thị vệ tinh từ IGS hoặc CelesTrak; khi mất mạng hoặc không có quỹ đạo phù hợp, giao diện vẫn hiển thị vị trí xấp xỉ thay vì bỏ trống. Việc gắn nhãn nguồn cũng làm cho kết quả minh bạch hơn với người dùng và thuận tiện khi gỡ lỗi.

\begin{table}[H]
\centering
\begin{tabular}{|p{2.2cm}|p{4.3cm}|p{6.4cm}|}
\hline
\textbf{Ưu tiên} & \textbf{Nguồn dữ liệu} & \textbf{Vai trò} \\ \hline
1 & Real GNSS PVT & Sử dụng tọa độ và vận tốc ECEF từ thiết bị nếu có thể trích xuất. \\ \hline
2 & IGS Broadcast Ephemeris & Lan truyền quỹ đạo từ lịch vệ tinh công khai khi PVT không có. \\ \hline
3 & CelesTrak TLE/SGP4 & Dùng phần tử quỹ đạo hai dòng cho GPS, Galileo, BeiDou và các chòm hỗ trợ. \\ \hline
4 & Approximate & Tạo vị trí xấp xỉ từ vị trí người quan sát, azimuth, elevation và bán kính quỹ đạo. \\ \hline
\end{tabular}
\caption{Chuỗi ưu tiên nguồn dữ liệu vệ tinh}
\label{table:gnss-source-priority}
\end{table}

\section{Chỉ đường trực tiếp bằng đo chuyển động thị giác khi GNSS suy giảm}
\label{section:5.live-routing-fusion}
\subsection{Bài toán}
Khi di chuyển bằng xe, GNSS có thể mất cập nhật trong hầm, khu đô thị dày đặc hoặc vùng bị che khuất. Nếu chỉ ngoại suy vị trí từ tốc độ cuối cùng trước khi mất tín hiệu, quỹ đạo sẽ sai nhanh mỗi khi xe tăng tốc, giảm tốc hoặc dừng. Ngược lại, nếu chỉ suy vận tốc từ độ lớn Optical Flow trung bình, kết quả cũng dễ sai vì chuyển động ảnh phụ thuộc vào ánh sáng, kết cấu mặt đường, rung camera, vật thể chuyển động trong cảnh và góc đặt điện thoại.

Bài toán đặt ra là một cơ chế hỗ trợ định vị ngắn hạn khi GNSS suy giảm, tận dụng camera, IMU và tuyến đường sẵn có, nhưng vẫn phát hiện được khi người dùng đi sai đường thay vì ép marker bám cứng vào tuyến cũ --- và không tự nhận là một hệ visual odometry đầy đủ.

\subsection{Giải pháp}
Đồ án triển khai một lớp điều phối chỉ đường trực tiếp hợp nhất dữ liệu theo hai chế độ. Khi GNSS còn tin cậy, vị trí và tốc độ từ dịch vụ định vị Android là nguồn chính, đồng thời hệ thống hiệu chỉnh hệ số quy đổi từ pixel/giây của Optical Flow sang mét/giây. Khi GNSS mất tin cậy, hệ thống chuyển sang ước lượng bằng dead reckoning: các chỉ số Optical Flow được gom thành một bộ đo chuyển động thị giác gồm tốc độ ước lượng, thay đổi hướng, độ dịch chuyển ảnh và hệ số chất lượng. Tốc độ xe trong giai đoạn hỗ trợ được trộn giữa tốc độ thị giác và tiền nghiệm từ tốc độ GNSS cuối cùng; tiền nghiệm này giảm dần theo thời gian mất tín hiệu nên hệ thống không giả định xe luôn giữ nguyên vận tốc, còn trọng số của camera tự giảm khi dòng chuyển động kém tin cậy.

Sau bước dự đoán vị trí, hệ thống thực hiện bám tuyến có điều kiện (chế độ SNAP) thay vì khớp bản đồ liên tục. Mỗi chu kỳ điều hướng khi GNSS mất tin cậy gồm các bước:
\begin{enumerate}
    \item Tạo bộ đo chuyển động thị giác từ độ dịch chuyển ngang/dọc, độ lớn trung bình và độ tin cậy của Optical Flow cùng tốc độ quay (yaw rate) của IMU.
    \item Tính tốc độ xe bằng cách trộn tốc độ thị giác với tiền nghiệm GNSS đang suy giảm, rồi giới hạn theo các ngưỡng tăng tốc, phanh và trôi theo quán tính.
    \item Dự đoán vị trí mới từ tốc độ và hướng hiện tại; độ bất định vị trí tăng dần theo thời gian mất GNSS và chất lượng camera.
    \item Chiếu vị trí dự đoán lên tuyến để tính độ tin cậy bám tuyến từ khoảng cách ngang, sai khác hướng, tính liên tục so với phép chiếu trước và chất lượng đo thị giác.
    \item Nếu độ tin cậy đủ cao, kéo một phần vị trí và hướng về tuyến rồi giảm độ bất định; nếu thấp, giữ nguyên vị trí dự đoán để không che giấu khả năng lệch tuyến.
\end{enumerate}
Khi quỹ đạo lệch tuyến qua nhiều mẫu liên tiếp và còn Internet, hệ thống yêu cầu một tuyến mới từ vị trí ước lượng tới đích.

\subsection{Kết quả đạt được}
Cơ chế này giúp chỉ đường trực tiếp phản ứng tốt hơn trong các khoảng mất GNSS ngắn, vì tốc độ hỗ trợ tăng giảm theo chuyển động ảnh thay vì cố định ở tốc độ cuối cùng. Bám tuyến có điều kiện làm giảm hiện tượng đoạn hỗ trợ trôi khỏi tuyến khi camera còn đáng tin, nhưng vẫn cho phép phát hiện lệch tuyến khi hướng và khoảng cách không còn phù hợp.

Tuy vậy, giải pháp vẫn có giới hạn rõ ràng: đây là cơ chế ước lượng chuyển động thực dụng dựa trên các chỉ số Optical Flow, chưa phải hệ thống khôi phục pose camera đầy đủ. Camera đơn không tự cung cấp tỷ lệ tuyệt đối nếu thiếu hiệu chỉnh bằng GNSS, và kết quả vẫn phụ thuộc vào ánh sáng, kết cấu bề mặt và điều kiện đặt thiết bị. Vì vậy, đúng như đánh giá ở Chương~\ref{chapter:Experiment}, chức năng này hiện ở mức thử nghiệm kỹ thuật và vẫn đang được tinh chỉnh.

\section{Cơ chế phân tích hiệu năng hai thuật toán Optical Flow}
\label{section:5.2}
\subsection{Bài toán}
KLT và Farneback có bản chất khác nhau --- KLT theo dõi các điểm đặc trưng thưa, còn Farneback ước lượng trường chuyển động dày đặc --- nên so sánh trực tiếp dễ sai lệch nếu hai thuật toán chạy ở thời điểm khác nhau hoặc dùng tham số không tương đương. Chỉ nhìn ảnh phủ vector, người dùng khó biết thuật toán nào nhanh hơn, ổn định hơn hay phát hiện nhiều chuyển động hơn trong cùng điều kiện.

Một yêu cầu khác là không trộn lẫn chỉ số đo trực tiếp với chỉ số mang tính suy luận. FPS và thời gian xử lý đo trực tiếp từ thời gian chạy, còn độ tin cậy lại phụ thuộc vào cách định nghĩa điểm khớp (inlier); nếu trình bày tất cả như các biểu đồ ngang hàng, người đọc dễ hiểu sai ý nghĩa từng chỉ số. Do đó cần phân tách chỉ số chính và chỉ số hỗ trợ.

\subsection{Giải pháp}
Đồ án triển khai một chế độ phân tích trên màn hình Optical Flow camera. Khi bật chế độ này, mỗi khung hình được nhân thành hai bản chạy song song qua KLT và Farneback dưới cùng điều kiện đầu vào, với độ nhạy hai thuật toán được khóa ở mức 100 để loại bỏ tác động của thao tác người dùng. Hai ảnh kết quả được ghép thành một khung trái--phải, mỗi nửa phủ thông số của thuật toán tương ứng để người dùng quan sát đồng thời.

Mỗi thuật toán trả về một bộ chỉ số riêng: KLT ghi số điểm theo dõi, số vector hoạt động, độ dịch chuyển trung vị và độ tin cậy theo sai số tiến-lùi (FBE); Farneback ghi số mẫu trên lưới, số vector vượt ngưỡng, độ dịch chuyển trung bình và độ tin cậy theo kiểm tra dòng chảy ngược. Để tách bạch hai loại chỉ số, phiên đo chỉ vẽ biểu đồ chính cho FPS, thời gian xử lý và số điểm/vector hoạt động (các đại lượng đo trực tiếp), còn độ tin cậy chỉ đưa vào phần tóm tắt để hỗ trợ diễn giải. Mỗi khung hình tạo ra một mẫu đo và được lưu vào phiên đang chạy.

\subsection{Kết quả đạt được}
Cơ chế tạo ra một quy trình đánh giá có tính lặp lại hơn so với quan sát thủ công từng thuật toán: người dùng chạy cùng một cảnh, dừng phiên rồi mở biểu đồ để xem diễn biến theo thời gian. Việc lưu phiên trong cơ sở dữ liệu Room giúp kết quả không phụ thuộc vòng đời màn hình và có thể mở lại sau. Bảng~\ref{table:analytics-metrics} mô tả ý nghĩa các chỉ số chính.

\begin{table}[H]
\centering
\begin{tabular}{|p{3.2cm}|p{4.4cm}|p{5.6cm}|}
\hline
\textbf{Chỉ số} & \textbf{Cách đo} & \textbf{Ý nghĩa} \\ \hline
FPS & $1000$ chia cho thời gian xử lý (ms) & Tốc độ xử lý tức thời của thuật toán trên khung hình hiện tại. \\ \hline
Thời gian xử lý & Thời gian chạy thuật toán đo bằng đồng hồ hệ thống & Chi phí CPU/GPU tương đối của từng phương pháp. \\ \hline
Số điểm/vector hoạt động & Số điểm KLT hoặc số mẫu Farneback vượt ngưỡng & Mức độ hoạt động của trường chuyển động được phát hiện. \\ \hline
Độ dịch chuyển trung bình & Trung bình hoặc trung vị của vector dịch chuyển & Hướng và độ lớn chuyển động biểu kiến sau lọc. \\ \hline
Độ tin cậy & Tỷ lệ điểm khớp theo kiểm tra tiến-lùi & Mức tin cậy tương đối, phụ thuộc vào định nghĩa sai số tiến-lùi của hệ thống. \\ \hline
\end{tabular}
\caption{Các thông số trong phiên phân tích Optical Flow}
\label{table:analytics-metrics}
\end{table}

\section{Quy trình tải lên và xử lý video RAFT trên máy chủ}
\label{section:5.3}
\subsection{Bài toán}
Mô hình RAFT cho trường Optical Flow dày đặc nhưng tốn nhiều bộ nhớ và thời gian suy luận, nên video được tải lên máy chủ để xử lý. Nếu gửi video bằng một yêu cầu đồng bộ duy nhất, kết nối dễ quá thời gian và người dùng không biết tiến độ; video lại là tệp lớn nên cần kiểm soát tệp tạm để không rò rỉ tài nguyên. Ngoài ra, người dùng không nên bị buộc ở lại màn hình gửi video: khi xử lý kéo dài, họ có thể chuyển màn hình hoặc khóa máy, do đó ứng dụng cần một cơ chế xử lý nền và một lớp phủ tiến độ độc lập với màn hình hiện tại.

\subsection{Giải pháp}
Đồ án tách xử lý RAFT khỏi thiết bị và đưa sang máy chủ, cho phép người dùng chọn backend tự triển khai hoặc backend ngoài. Với backend tự triển khai, video nhỏ được gửi trực tiếp qua yêu cầu tạo job, còn video lớn được tải theo phiên nhiều phần rồi gọi yêu cầu hoàn tất để máy chủ tạo job; sau khi nhận mã job, ứng dụng thăm dò trạng thái theo chu kỳ (đang chờ, đang xử lý, hoàn tất, thất bại, đang hủy, đã hủy). Tác vụ xử lý chạy nền, cập nhật tiến độ qua một hàm gọi lại và kiểm tra cờ hủy.

Trong bộ xử lý của máy chủ, mỗi video được đọc bằng OpenCV với một vùng đệm khung hình ngắn để ước lượng chuyển động qua khoảng cách vài khung hình, giúp vector rõ hơn ở cảnh chuyển động chậm. Kết quả được vẽ theo hai chế độ: chế độ vector lấy mẫu theo lưới, lọc vector yếu theo phân vị và chuẩn hóa chiều dài hiển thị; chế độ heatmap tính độ lớn chuyển động, làm mượt, chuẩn hóa rồi trộn bảng màu với khung hình gốc. Video kết quả được ghi bằng ffmpeg theo chuẩn H.264 để Android phát ổn định. Tệp tạm đầu vào được xóa sau khi job kết thúc; kết quả thành công được giữ tới khi ứng dụng tải về hoặc yêu cầu dọn dẹp, còn kết quả lỗi hoặc bị hủy được xóa ngay.

Khi người dùng chọn vùng đối tượng trên video, ứng dụng gửi kèm tọa độ vùng chọn, tỷ lệ khung xem trước và mốc thời gian được chọn để máy chủ tạo mặt nạ khởi tạo tại đúng khung hình. Nếu dùng Cutie, mặt nạ được lan truyền tiến và lùi theo thời gian để bám đối tượng qua toàn bộ video, và Optical Flow chỉ được vẽ trong vùng mặt nạ đang hoạt động; nhờ đó đối tượng vẫn được theo dõi ngay cả khi rời khỏi khung chữ nhật ban đầu. Khác với cơ chế bám đối tượng trên camera trực tiếp ở Mục~\ref{section:5.4}, hướng xử lý này có toàn bộ video nên lan truyền được mặt nạ hai chiều thay vì chỉ bám theo thời gian thực.

Ở phía Android, một tác vụ nền chạy ở chế độ foreground qua WorkManager chịu trách nhiệm tải video lên, thăm dò trạng thái, tải kết quả về và hủy job khi backend hỗ trợ, nhờ vậy hệ thống không dừng tác vụ dài khi người dùng rời màn hình. Lớp phủ tiến độ được quản lý ở cấp Activity, có thể hiển thị dạng thẻ hoặc thu gọn thành bong bóng nổi, và chuyển sang trạng thái hoàn tất khi video sẵn sàng.

\subsection{Kết quả đạt được}
Quy trình bất đồng bộ này mang lại trải nghiệm rõ ràng hơn khi xử lý video dài: ứng dụng không bị khóa ở một màn hình cụ thể, người dùng có thể mở video kết quả ngay khi job thành công hoặc nhận thông báo rõ ràng khi job lỗi hay bị hủy. Về phía máy chủ, việc tách tài nguyên tạm đầu vào/đầu ra, giới hạn số job đồng thời, dọn dẹp chủ động và tải kết quả theo từng phần giúp giảm nguy cơ quá tải và tích tụ tệp tạm sau nhiều lần thử nghiệm.

\section{Cơ chế bám đối tượng trong luồng Optical Flow camera trực tiếp}
\label{section:5.4}
\subsection{Bài toán}
Nhiều khi người dùng chỉ quan tâm chuyển động của một đối tượng cụ thể --- một phương tiện, một người đi bộ --- thay vì toàn bộ chuyển động nền. Xử lý toàn khung hình vừa dễ nhiễu bởi vùng không liên quan, vừa làm tăng thời gian xử lý. Cho người dùng khoanh một vùng cố định là cách đơn giản, nhưng trên camera trực tiếp cả đối tượng lẫn thiết bị đều liên tục dịch chuyển, nên một hình chữ nhật tĩnh sẽ lệch khỏi đối tượng chỉ sau vài khung hình.

Bài toán là một cơ chế bám đối tượng hoạt động ngay trên luồng camera thời gian thực: người dùng chọn đối tượng một lần, hệ thống tự cập nhật vùng phân tích theo đối tượng ở các khung hình sau. Cơ chế này phải đủ nhẹ để giữ yêu cầu thời gian thực và đủ ổn định để không trôi sang nền khi đối tượng đổi kích thước hoặc bị che một phần.

\subsection{Giải pháp}
Đồ án triển khai một lớp phủ để người dùng khoanh đối tượng cần theo dõi; vùng chọn được lưu dưới dạng tọa độ chuẩn hóa theo khung nhìn. Ở mỗi khung hình, ứng dụng ánh xạ tọa độ chuẩn hóa về tọa độ khung hình, chạy Optical Flow trong vùng con và khôi phục phần ngoài vùng bằng khung hình gốc để giữ nguyên bối cảnh hiển thị. Điểm cốt lõi là vùng con này không cố định mà được cập nhật bởi một bộ bám đối tượng: hệ thống dùng bộ theo dõi CSRT của OpenCV~\cite{opencv-video-tracking} kết hợp đối sánh mẫu (template matching). CSRT cho chất lượng bám tốt với đối tượng biến dạng nhẹ, còn template matching đóng vai trò kiểm tra và khôi phục khi tracker mất dấu trong thời gian ngắn.

Khó khăn chính là khung xem trước của camera có thể bị co giãn và cắt khác với khung hình xử lý, nên phép ánh xạ phải tính tỷ lệ khung nhìn, kích thước khung hình, hệ số co giãn và độ lệch thì hộp bám mới khớp với những gì người dùng nhìn thấy. Nếu vùng chọn quá nhỏ, hệ thống yêu cầu chọn lại vì Optical Flow cần đủ pixel và bộ bám cần đủ kết cấu để bám ổn định. Khi tracker báo độ tin cậy thấp hoặc kết quả lệch quá xa vị trí dự đoán, hệ thống giữ lại hộp bám gần nhất thay vì nhảy đột ngột, tránh hiện tượng vùng phân tích rung hoặc trôi sang nền. Khác với cơ chế bám cho video ngoại tuyến ở Mục~\ref{section:5.3} --- nơi máy chủ có toàn bộ chuỗi khung hình nên dùng Cutie lan truyền mặt nạ hai chiều --- trên camera trực tiếp hệ thống chỉ có khung hình quá khứ và hiện tại nên phải bám theo thời gian thực bằng CSRT.

\subsection{Kết quả đạt được}
Bộ bám giúp vùng phân tích đi theo đối tượng trong thời gian thực thay vì cố định ở hình chữ nhật ban đầu, nhờ đó người dùng tập trung quan sát một mục tiêu, giảm nhiễu nền và giảm chi phí xử lý do chỉ chạy Optical Flow trong vùng con. Chức năng này cũng làm ứng dụng giống một công cụ thử nghiệm hơn vì có thể chọn lại đối tượng cần theo dõi mà không cần sửa mã nguồn. Hạn chế còn lại là bộ bám vẫn có thể mất dấu khi đối tượng bị che hoàn toàn, ra khỏi khung hình hoặc khi cảnh thiếu kết cấu; trong các trường hợp đó người dùng cần khoanh lại đối tượng để khôi phục việc bám.

\section{Bộ dựng giao diện kính lỏng thời gian thực bằng AGSL với cơ chế suy biến mượt}
\label{section:5.liquidglass}
\subsection{Bài toán}
Giao diện ``kính lỏng'' (liquid glass) là lớp kính trong mờ có khúc xạ, làm mờ nền phía sau và viền ánh sáng tán sắc, tạo cảm giác hiện đại và có chiều sâu. Trong đồ án, hiệu ứng này được áp dụng cho các bảng điều khiển nổi trên màn hình bản đồ. Đặc thù của màn hình này là khi người dùng kéo hoặc phóng to bản đồ, nền nằm dưới lớp kính thay đổi liên tục, nên hiệu ứng phải được tính lại ở mỗi khung hình mà không gây giật hình. Đây cũng là phần tốn kém nhất, bởi việc khúc xạ, làm mờ và tán sắc trên từng điểm ảnh là rất nặng nếu thực hiện bằng cách vẽ thông thường trên bộ xử lý trung tâm.

Khó khăn thứ hai đến từ sự phân mảnh của hệ điều hành Android. Những công cụ tô bóng đồ họa mạnh nhất chỉ xuất hiện trên các phiên bản mới: khả năng chạy shader đồ họa thời gian thực bằng ngôn ngữ AGSL chỉ có từ Android 13~\cite{android-agsl}, còn khả năng làm mờ bằng phần cứng chỉ có từ Android 12~\cite{android-rendereffect}; trên những thiết bị cũ hơn, các khả năng này không tồn tại. Bài toán đặt ra là dựng được hiệu ứng kính lỏng chất lượng cao trên thiết bị đời mới nhưng vẫn suy biến mượt mà trên thiết bị đời cũ --- giữ được vẻ ngoài ``kính'' mà không khiến ứng dụng dừng đột ngột vì gọi tới khả năng thiết bị không hỗ trợ.

\subsection{Giải pháp}
Giải pháp là một bộ dựng kính lỏng được thiết kế theo nhiều mức chất lượng. Khi khởi tạo, bộ dựng kiểm tra phiên bản Android của thiết bị và tự chọn mức cao nhất mà thiết bị có thể chạy. Toàn bộ việc lựa chọn này được che giấu phía sau một giao diện chung, nên màn hình bản đồ chỉ cần yêu cầu ``một lớp kính'' mà không phải bận tâm bên dưới đang dùng cách dựng nào. Ba mức được tổ chức từ cao xuống thấp và được tóm tắt trong Bảng~\ref{table:liquidglass-tiers}.

Ở mức cao nhất, dành cho thiết bị từ Android 13 trở lên, hệ thống chụp lại phần bản đồ nằm dưới lớp kính rồi đưa qua một shader đồ họa chạy trực tiếp trên bộ xử lý đồ họa. Shader này tạo ra hiệu ứng khúc xạ uốn theo độ bo góc và độ dày mép kính, cảm giác chiều sâu, hiện tượng tán sắc màu ở rìa, cùng với việc điều chỉnh tương phản, độ sáng, độ bão hòa và lớp màu phủ; ngoài ra còn có thể chồng thêm một lớp làm mờ. Vì mọi tính toán nặng đều dồn cho bộ xử lý đồ họa nên hiệu ứng cập nhật được theo từng khung hình ở đúng tốc độ làm tươi của màn hình.

Ở mức trung gian, dành cho Android 12, thiết bị chưa chạy được shader nói trên nhưng vẫn làm mờ được bằng phần cứng. Hệ thống tận dụng khả năng làm mờ này cho phần nền, rồi vẽ chồng lên một lớp kính mô phỏng: các dải sáng chuyển màu gợi cảm giác phản chiếu và chiều sâu, cùng một viền hai màu lệch nhau để tái hiện hiệu ứng tán sắc mà mức cao nhất tạo ra bằng shader.

Ở mức thấp nhất, dành cho các thiết bị cũ hơn Android 12 vốn không còn khả năng làm mờ bằng phần cứng, hệ thống chụp nền vào một ảnh thu nhỏ, làm mờ ảnh này bằng thuật toán chạy trên bộ xử lý trung tâm với tần suất cập nhật được giới hạn để tránh quá tải, rồi vẽ chồng cùng lớp kính mô phỏng như ở mức trung gian. Việc thu nhỏ ảnh trước khi làm mờ giúp giảm đáng kể khối lượng tính toán trong khi mắt thường khó nhận ra khác biệt.

Bên cạnh việc tự chọn mức, bộ dựng còn được tối ưu để hầu như không tạo thêm đối tượng mới trong mỗi khung hình: kết quả làm mờ được tái sử dụng và chỉ làm tươi theo chu kỳ, còn lớp kính chỉ được dựng lại khi một tham số thật sự thay đổi. Hệ thống cũng theo dõi thời điểm bản đồ vẽ lại để luôn chụp đúng nội dung mới, đồng thời cung cấp một tùy chọn ép chạy một mức bất kỳ nhằm thử nghiệm hiệu ứng trên nhiều dòng thiết bị khác nhau.

\subsection{Kết quả đạt được}
Cùng một thành phần giao diện cho ra hiệu ứng kính khúc xạ đầy đủ trên thiết bị từ Android 13 trở lên, kính mờ bằng phần cứng trên Android 12 và kính mờ dựng trên bộ xử lý trung tâm với thiết bị cũ hơn, mà màn hình bản đồ không cần thay đổi cách sử dụng. Nhờ tự chọn mức theo phiên bản hệ điều hành, ứng dụng không bị dừng đột ngột vì gọi tới khả năng không được hỗ trợ, đồng thời vẫn giữ được vẻ ngoài ``kính'' trên mọi thiết bị. Ở mức cao nhất, do phần lớn khối lượng tính toán nằm trên bộ xử lý đồ họa và được tái sử dụng, hiệu ứng vẫn mượt ngay cả khi người dùng kéo và phóng bản đồ liên tục. Đây là một ví dụ cụ thể cho chiến lược tăng cường lũy tiến (progressive enhancement): khai thác tối đa năng lực của thiết bị mới nhưng vẫn bảo đảm trải nghiệm chấp nhận được trên thiết bị cũ.

\begin{table}[H]
\centering
\begin{tabular}{|p{2.7cm}|p{4.0cm}|p{6.2cm}|}
\hline
\textbf{Phiên bản Android} & \textbf{Kỹ thuật} & \textbf{Mức hiệu ứng} \\ \hline
13 trở lên & Shader đồ họa AGSL chạy trên bộ xử lý đồ họa & Khúc xạ, tán sắc, chiều sâu, tương phản, màu phủ và làm mờ đầy đủ. \\ \hline
12 & Làm mờ bằng phần cứng kết hợp lớp kính mô phỏng & Làm mờ bằng phần cứng; mô phỏng khúc xạ và tán sắc bằng dải sáng chuyển màu và viền lệch màu. \\ \hline
Trước 12 & Làm mờ ảnh thu nhỏ trên bộ xử lý trung tâm kết hợp lớp kính mô phỏng & Làm mờ cơ bản; giữ vẻ ngoài kính nhưng không có khúc xạ thật. \\ \hline
\end{tabular}
\caption{Ba mức hiện thực hiệu ứng kính lỏng theo phiên bản Android của thiết bị}
\label{table:liquidglass-tiers}
\end{table}

\section*{Kết chương}
Chương này đã trình bày sáu giải pháp mang tính đóng góp của đồ án: chuỗi phân giải vị trí vệ tinh nhiều nguồn, cơ chế hỗ trợ định vị khi GNSS suy giảm, phương pháp so sánh hiệu năng giữa hai thuật toán Optical Flow, quy trình xử lý video RAFT bất đồng bộ trên máy chủ, cơ chế bám đối tượng trên camera trực tiếp và bộ dựng giao diện kính lỏng thời gian thực bằng AGSL kèm cơ chế suy biến mượt. Mỗi giải pháp đều xuất phát từ một hạn chế thực tế và được kiểm chứng qua kết quả chạy thử trên hệ thống. Những kết quả và đóng góp này, cùng các hạn chế còn lại, được tổng kết trong Chương~\ref{chapter:conclusion}.

\end{document}
